\section*{Chapter \ref{ch:b3:computability} --- Turing Machines and Computability}

\begin{solution}{exo:b3:computability:1}
\texttt{[q0]aba}, \texttt{\_[has\_a]ba}, \texttt{\_b[has\_a]a}, \texttt{\_ba[has\_a]\_}, \texttt{\_b[end\_a]a\_}, \texttt{\_[back]b\_\_}. The machine has erased the first \texttt{a}, walked to the end, found an \texttt{a} there, erased it, and is walking back. Five more steps accept: the \texttt{b} left alone is erased and the machine finds blanks on both sides.
\end{solution}

\begin{solution}{exo:b3:computability:2}
(a)~\omterm{def:b3:computability:languages}{Decidable}: read the \omterm{def:b3:computability:universal}{encoding} and count the states. (b)~\omterm{def:b3:computability:languages}{Recognisable}, since a machine can run $N$ on the empty word through the universal machine and accept if it accepts; not \omterm{def:b3:computability:languages}{decidable}, by the reduction of the chapter. (c)~Not \omterm{def:b3:computability:languages}{recognisable}: if it were, $E$ and its complement would both be \omterm{def:b3:computability:languages}{recognisable}, and $E$ would be \omterm{def:b3:computability:languages}{decidable} by \cref{thm:b3:computability:complement}. (d)~\omterm{def:b3:computability:languages}{Decidable}: run $M$ on $w$ for at most 1,000 steps.
\end{solution}

\begin{solution}{exo:b3:computability:3}
$(\mathtt{q0}, \mathtt{+}) \to (\mathtt{q0}, \mathtt{1}, R)$ is the numbers 1, 1, 1, 2, 2: \texttt{10101011011}. $(\mathtt{last}, \mathtt{1}) \to (\mathtt{done}, \mathtt{\_}, S)$ is 3, 2, 2, 3, 3: \texttt{11101101101110111}. The whole adder has 58 characters (\code{utm.encode(ADD)}).
\end{solution}

\begin{solution}{exo:b3:computability:4}
The increment walks right over four digits (4 steps), steps back onto the blank's left (1), turns four 1s into 0s (4) and writes a 1 on the blank (1): tape \texttt{10000}, 10 steps. The adder on \texttt{1+1} turns \texttt{+} into 1 and erases the last 1: \texttt{11}. The palindrome machine on \texttt{aa} erases the first \texttt{a}, passes the second, finds the blank, steps back, erases the second, steps back, and finds an empty word: 6 steps, $(2 + 1)(2 + 2)/2$ (\code{test_exercise_4_predictions}).
\end{solution}

\begin{solution}{exo:b3:computability:5}
On \texttt{111} the carry turns all three 1s into 0s and reaches the blank left of the number, where no transition applies: the machine halts with \texttt{000}. The missing transition is \texttt{("carry", B):\ ("done", "1", "S")}, a new leading 1 (\code{exercises/ex5_increment.py}).
\end{solution}

\begin{solution}{exo:b3:computability:6}
\code{exercises/ex6_anbn.py}: each round crosses off the leftmost \texttt{a} as \texttt{X}, walks right over \texttt{a}s and \texttt{Y}s to the first \texttt{b}, crosses it off as \texttt{Y}, and walks back to the last \texttt{X}. When the next letter is a \texttt{Y} instead of an \texttt{a}, every \texttt{a} is crossed off, and the machine accepts if only \texttt{Y}s follow. A \texttt{b} before an \texttt{a}, an \texttt{a} after the \texttt{b}s, or a missing \texttt{b} leaves the machine with no transition, and it rejects.
\end{solution}

\begin{solution}{exo:b3:computability:7}
Let $R_1$ recognise $L$ and $R_2$ recognise $\overline{L}$. The \omterm{def:b3:computability:languages}{decider} $D$, on $w$, simulates one step of $R_1$ on $w$, then one step of $R_2$, and so on, keeping both \omterm{def:b3:computability:tm}{configurations} on its tape (a two-tape machine, then one tape by \cref{thm:b3:computability:multitape}). If $R_1$ accepts, $D$ accepts; if $R_2$ accepts, $D$ rejects. Every word is in $L$ or in $\overline{L}$, so one of the two accepts after finitely many steps, and $D$ halts. It never answers wrongly, since $R_1$ accepts only words of $L$ and $R_2$ only words of $\overline{L}$.
\end{solution}

\begin{solution}{exo:b3:computability:8}
The map $\langle M, w \rangle \mapsto \langle M_w \rangle$ of the chapter works: $M_w$ ignores its input and runs $M$ on $w$, so it halts on every input if $M$ halts on $w$ and on none otherwise. Hence $\mathit{HALT} \le_m \mathit{TOTAL}$, and $\mathit{TOTAL}$ is undecidable. It is not a case of Rice's theorem as stated: halting on every input is not a property of the language accepted, since a machine that rejects by halting and one that rejects by looping can accept the same words.
\end{solution}

\begin{solution}{exo:b3:computability:9}
Two tapes: $n$ steps copy, one step at the end of the word, $n$ steps rewind and one more step back onto the first letter, then $n$ comparisons and one step to accept: $3n + 3$. One tape: each round on a word of length $m$ walks right over $m$ cells and back over $m - 1$, plus the steps onto the ends, about $2m$ steps, and removes two letters; summing over $m = n, n - 2, \ldots$ gives $(n + 1)(n + 2)/2$, quadratic as \cref{thm:b3:computability:multitape} allows (\code{test_palindromes_one_and_two_tapes} checks both formulas for every palindrome of length at most 8).
\end{solution}

\begin{solution}{exo:b3:computability:10}
If the machines with empty language have $\Pi$, then $\overline{\Pi}$, the property of lacking $\Pi$, is \omterm{def:b3:computability:rice}{semantic} and \omterm{def:b3:computability:rice}{non-trivial}, and machines with the empty language lack it. By the case proved in the chapter, $\overline{\Pi}$ is undecidable. A \omterm{def:b3:computability:languages}{decider} for $\Pi$, with its answers swapped, would decide $\overline{\Pi}$; so $\Pi$ is undecidable too.
\end{solution}

\begin{solution}{exo:b3:computability:11}
Facts used: $f \ge 0$ is total and computable, so $g(n) = f(2n) + f(2n + 1) + 1$ is too, and a two-symbol machine computes it from $n$ written in unary, with some number $c$ of states, leaving $g(n)$ 1s; writing $n$ 1s on the blank tape takes $n$ states, one per 1; $S$ is increasing (problem, question 18).

The machine $B_n$ writes $n$ 1s, then runs the machine for $g$. It halts, has $n + c$ states, and takes at least $g(n)$ steps, since it leaves $g(n)$ 1s and writes one per step at most. So $S(n + c) \ge g(n)$. For $n \ge c$, $n + c \le 2n$, hence $S(2n + 1) \ge S(2n) \ge S(n + c) \ge g(n)$, which exceeds both $f(2n)$ and $f(2n + 1)$. So $S(m) > f(m)$ for every $m \ge 2c$. In particular $S \ne f$ for every computable $f$: a second proof of \cref{thm:b3:computability:busy-beaver}.
\end{solution}

\begin{solution}{exo:b3:computability:12}
\code{exercises/ex12_primrec.py}. \code{pred} is primitive recursion with base 0 and step ``the previous $n$''. \code{monus} recurses on the subtracted number, applying \code{pred} once per step, then swaps the arguments. A test for zero, $1 \mathbin{\dot-} x$, and equality, zero of $(x \mathbin{\dot-} y) + (y \mathbin{\dot-} x)$, give the remainder: count up and wrap to 0 when the count reaches $y$. \code{div} adds 1 each time the remainder of the next number is 0 (\code{test_exercise_12_primitive_recursive}).
\end{solution}

\begin{solution}{pb:b3:computability:1}
\textbf{1.} Write 0 or 1, move left or right, go to one of $n + 1$ states: $4(n + 1)$ actions. Each of the $2n$ pairs (state, symbol) gets one: $(4(n + 1))^{2n}$ machines.

\textbf{2.} $8^2 = 64$ and $12^4 = 20{,}736$. For three states, $16^6 = 16{,}777{,}216$.

\textbf{3.} From the blank tape the only transition ever used is A0. Of its 8 actions, 4 go to H and 4 back to A: those 4 times the 8 choices for the unused A1 make 32 halting machines, each after one step. The others repeat A0 for ever.

\textbf{4.} $24^{10} \approx 6.3 \times 10^{13}$. A proof needs one representative per class of equivalent machines: it fixes the first move, ignores mirror images, and builds the table only as far as the run uses it, filling a transition the first time it is reached; unused transitions never multiply the count.

\textbf{5.} Every run is the mirror image of the other's: the tape is reversed, and the steps and the 1s are the same.

\textbf{6.} Some machines never halt, and no program can tell in general which: without a cap the enumeration would never end.

\textbf{7.} It can count as non-halting a machine that halts after the cap, and so underestimate $S$ and $\Sigma$; it never overestimates them.

\textbf{8.} Every halting two-state machine halts within $S(2) = 6$ steps, so one that is still running at step 100 never halts.

\textbf{9.} Only $S(2) \ge 6$: some machine might halt after step 100.

\textbf{10.} A machine could halt after step 1,000; a larger cap pushes the doubt back without removing it. The proof of $S(2) = 6$ shows for each remaining machine why it never halts.

\textbf{11.} From the blank tape, consider the states the machine can reach by following the transitions from A. If none of them has a transition to H, then at every step the machine is in one of these states, and some transition applies (every pair has one), so it never halts.

\textbf{12.} \texttt{[A]0}, \texttt{[B]01}, \texttt{1[A]1}, \texttt{11[B]0}, \texttt{111[A]0}, \texttt{11[B]11}, \texttt{1[H]111}, with 0 for blank: 6 steps, four 1s.

\textbf{13.} The halting transition is B1: it reads a 1. Writing 1 keeps it; writing 0 would erase it and leave three 1s.

\textbf{14.} \texttt{1RB 1LB 1LA 1RH}: every move reversed, the same run reflected.

\textbf{15.} The champion and its mirror, each with the halting move left or right: the last move does not change the tape.

\textbf{16.} The steps do not depend on what the halting transition writes, nor on its move; and several different runs of six steps exist, leaving fewer 1s.

\textbf{17.} Yes: 40 machines take six steps, only 4 leave four 1s.

\textbf{18.} Take a machine with $n$ states that takes $S(n)$ steps. Replace its halting transition by one to a new state, which then halts in one more step: $n + 1$ states, $S(n) + 1$ steps.

\textbf{19.} Given $M$ with $n$ states, compute $S(n)$ and run $M$ for $S(n)$ steps from the blank tape. If it has not halted, it never will, since a halting machine of $n$ states takes at most $S(n)$ steps.

\textbf{20.} Halting on the blank tape is undecidable: $\mathit{HALT}$ reduces to it by $\langle M, w \rangle \mapsto \langle M_w \rangle$, $M_w$ written with two symbols.

\textbf{21.} If it were computable, it would exceed itself at all large $n$.

\textbf{22.} A single program giving them would compute $S$: from the proofs, enumerate the machines, discard those proved non-halting, and run the others to the end.

\textbf{23.} A computer counted the 181 million cases; the proof assistant checks that the program which did so, and the argument it implements, are correct, rather than asking readers to trust the code.

\textbf{24.} An enumeration with a cap gives lower bounds for $S(n)$ and $\Sigma(n)$, and a proof that the remaining machines never halt makes them exact: $S(2) = 6$ and $\Sigma(2) = 4$ by the chapter's program, $S(5) = 47{,}176{,}870$ by a collaboration. No program computes $S(n)$ for every $n$, since that program would decide halting, and $S$ outgrows every computable function.
\end{solution}
